\documentclass{article}
\usepackage{graphicx}
\usepackage{mathrsfs}
\usepackage{amsmath}
\usepackage{amssymb}

\title{Connectivity Matrices}
\author{Oneal Wang}
\date{October 2025}

\begin{document}
    \maketitle

    \section{Introduction}

    \section{Shale}
    $V$ is represented as a matrix of the sending node as rows and receiving node
    as columns
    \[
        V = \mathscr{M}_{n \times n}
    \]
    $D$ is the sending data represented as a matrix with b bits and n nodes
    \[
        D = \mathscr{M}_{b \times n}
    \]
    $T^{i}$ transformation from V to W for some timeslot i (cyclical) The full edge
    basis is trivial: $e_{i,j}$ is the standard unit vector with 1 at the
    coordinate for edge (i,j) and 0 elsewhere. Those $m=(nC2)$ vectors span $R^{m}$.
    The round-robin matching vectors are then simply sums of some of those edge
    unit vector s.
    \[
        T^{i} = \mathscr{M}_{n \times n}
    \]
    \[
        T^{i} \rightarrow\mathscr{L}(V,W)
    \]
    $W$ is the recieving node for each sending node for a specific time slot called
    an adjacency matrix
    \[
        W = \mathscr{M}_{n \times n}\mid W_{A_{i,k},i}= 1, \neg(W_{A_{i,k},i}) =
        0
    \]
    $S$ is the recieved data represented as a matrix with b bits and n nodes
    \[
        S = \mathscr{M}_{b \times n}
    \]
    \begin{enumerate}
        \item Round-Robin 1 Letter 1 Node\\
            minimum time coefficient: n\\
            maximum time coefficient: n(n-1)\\
            number of pathways: n(n-1)/2\\
            for an arbitrary i broken nodes, remove $(2ni-i+1)/2$ pathways\\
            for an arbitrary k broken pathways, maximum increase to n-1 timeslots,
            centered around +1 timeslot
            \[
                \mathscr{M}_{b \times n}=DT^{i}(V) = DW = S, i \in \{1, ... ,n-1\}
                \\
            \]
            Example 2.1 (6 nodes, 5 timeslots, h=1): \\
            W = Adjacency Matrix of a Directed Graph, where each timeslot is the
            row and each column is the node
            \[
                A =
                \begin{bmatrix}
                    2 & 3 & 4 & 5 & 6 & 1 \\
                    3 & 4 & 5 & 6 & 1 & 2 \\
                    4 & 5 & 6 & 1 & 2 & 3 \\
                    5 & 6 & 1 & 2 & 3 & 4 \\
                    6 & 1 & 2 & 3 & 4 & 5 \\
                \end{bmatrix}
            \]
            \[
                T =
                \begin{bmatrix}
                    0 & 0 & 0 & 0 & 0 & 1 \\
                    1 & 0 & 0 & 0 & 0 & 0 \\
                    0 & 1 & 0 & 0 & 0 & 0 \\
                    0 & 0 & 1 & 0 & 0 & 0 \\
                    0 & 0 & 0 & 1 & 0 & 0 \\
                    0 & 0 & 0 & 0 & 1 & 0 \\
                \end{bmatrix}
            \]
            \[
                V =
                \begin{bmatrix}
                    0 & 0 & 0 & 0 & 0 & 1 \\
                    1 & 0 & 0 & 0 & 0 & 0 \\
                    0 & 1 & 0 & 1 & 0 & 0 \\
                    0 & 0 & 1 & 1 & 0 & 0 \\
                    1 & 0 & 0 & 1 & 0 & 0 \\
                    0 & 0 & 0 & 0 & 1 & 0 \\
                \end{bmatrix}
            \]
            \[
                W_{5 \times 5}= \mathscr{M}\mid W_{A_{i,k},i}= 1, \neg(W_{A_{i,k},i}
                ) = 0
            \]
            \[
                DT^{i}(V_{6 \times 6}) = DW_{6 \times 6}= D
                \begin{bmatrix}
                    0 & 0 & 0 & 0 & 1 & 0 \\
                    0 & 0 & 0 & 0 & 0 & 1 \\
                    1 & 0 & 1 & 0 & 0 & 0 \\
                    0 & 1 & 1 & 0 & 0 & 0 \\
                    0 & 0 & 1 & 0 & 0 & 1 \\
                    0 & 0 & 0 & 1 & 0 & 0 \\
                \end{bmatrix}
                = S , i \in \{1, ... ,5\}
            \]
            minimum time coefficient: 6\\
            maximum time coefficient: 30\\
            number of pathways: 15\\

        \item Round-Robin 1 Letter repeat for every X Letters of 1 Node\\
            number of nodes: $L^{x}$ = n \\
            timeslot number: $(L-1)^{x}$ \\
            number of pathways: $L^{x} \times (L-1)^{X}/2$\\
            representation in base L for matrix algebra\\
            for each node k, available timeslots are exactly one digit
            difference from k in base L\\
            for an arbitrary i broken nodes, remove $i(L-1)^{x}$ pathways
            \[
                \mathscr{M}_{b \times n}= DW_{i}= S, i \in \{0, ... ,(L-1)^{x}\},
                j \in \mathbb{Z}, 0\leq j \leq L^{x}-1\

            \]
            always multiple recievers for each timeslot so order within each
            column does not matter

            Example 2.2 (3 letters, x = 2): \\
            number of nodes: $3^{2}$ = 9 \\
            timeslot number: $2^{2} = 4$ \\
            number of pathways: $3^{2} \times 2^{2}/2 = 18$\\
            W = Adjacency Matrix of a Directed Graph, where each timeslot is the
            row and each column is the node
            \[
                A =
                \begin{bmatrix}
                    1 & 0 & 0 & 0 & 1 & 2 & 0 & 1 & 2 \\
                    2 & 2 & 1 & 4 & 3 & 3 & 3 & 4 & 5 \\
                    3 & 4 & 5 & 5 & 5 & 4 & 7 & 6 & 6 \\
                    6 & 7 & 8 & 6 & 7 & 8 & 8 & 8 & 7 \\
                \end{bmatrix}
            \]
            \[
                W_{8 \times 8}= \mathscr{M}\mid W_{A_{i,k},i}= 1, \neg(W_{A_{i,k},i}
                ) = 0
            \]
            \[
                \mathscr{M}_{b \times n}= DW = S, i \in \{1, ... ,n-1\}\\
            \]

        \item Round-Robin t Letter repeat (Hamming-distance) for every X Letters
            of 1 Node\\
            number of nodes: $L^{x}$ = n \\
            timeslot number: $(L-1)^{x+t-1}$ \\
            number of pathways: $L^{x}(L^{x}-1)^{x+t-1}/2$ \\
            representation in base L for matrix algebra\\
            for each node k, available timeslots are exactly t digit difference from
            k in base, also the hamming distance L\\
            \[
                \mathscr{M}_{b \times n}= DW_{i}= S, i \in \{0, ... ,(L-1)^{x+t-1}
                \}, j \in \mathbb{Z}, 0\leq j \leq (L-1)^{x+t-1}\

            \]
            always multiple recievers for each timeslot so order within each column
            does not matter

            Example 2.3 (4 letters, x = 3, t=2): \\
            number of nodes: $4^{2}$ = 16 \\
            timeslot number: $3^{3} = 27$ \\
            number of pathways: $4^{2} \times 3^{3}/2 = 216$\\
            W = Adjacency Matrix of a Directed Graph, where each timeslot is the
            row and each column is the node
            \[
                A_{27 \times 27}=
                \begin{bmatrix}
                    5   & 4  & ... & 2  & 3   \\
                    6   &    &     &    & 7   \\
                    ... &    & ... &    & ... \\
                    56  &    &     &    & 57  \\
                    60  & 61 & ... & 59 & 58  \\
                \end{bmatrix}
            \]
            \[
                W_{27 \times 27}= \mathscr{M}\mid W_{A_{i,k},i}= 1, \neg(W_{A_{i,k},i}
                ) = 0
            \]
            \[
                \mathscr{M}_{b \times 27}= DW = S, i \in \{1, ... ,n-1\}\\
            \]
    \end{enumerate}
    \section{Opera}
    $V$ is represented as a matrix of the racks with the sending node as rows and
    receiving node as columns
    \[
        V = \mathscr{M}_{n \times n}
    \]
    $D$ is the sending data represented as a matrix with b bits and n nodes
    \[
        D = \mathscr{M}_{b \times n}
    \]
    \begin{enumerate}
        \item Direct Path \\
            timeslots = maixmum number of hops for an arbitrary node\\
            for an arbitrary i broken racks, maximum of n-1 hops\\
            Example 3.1 (8 ToR Switches, 4 Rotor Circuit Switches): \\
            A = represents a directed graph where each original row is the node
            and each column is the rack
            \[
                A =
                \begin{bmatrix}
                    3 & 7 & 1 & 8 & 5 & 2 & 6 & 4 \\
                    7 & 8 & 4 & 5 & 3 & 1 & 2 & 6 \\
                    1 & 4 & 5 & 3 & 2 & 6 & 7 & 8 \\
                    5 & 3 & 2 & 4 & 6 & 7 & 8 & 1 \\
                    4 & 6 & 3 & 2 & 1 & 8 & 1 & 7 \\
                    6 & 5 & 8 & 7 & 4 & 3 & 3 & 2 \\
                    2 & 1 & 7 & 6 & 8 & 4 & 4 & 5 \\
                    8 & 2 & 6 & 1 & 7 & 5 & 5 & 3 \\
                \end{bmatrix}
            \]
            W = Adjacency Matrix of a Directed Graph, where each timeslot is the
            row and each column is the node \\
            \[
                W_{i}= \mathscr{M}\mid W_{A_{i,k},i}= 1, \neg(W_{i_{A_{i,k},i}})
                = 0
            \]

        \item (Uniform) Indirect Path \\
            This uses BFS to find the optimal path from node a to node b. \\
            Minimum of 1 hop and maximum of
            $\left\lceil \frac{n}{2}\right\rceil$ hops\\
            For an arbitrary i broken racks, at least no hop up to +i hops change\\
            number of timeslots = maximum number of hops \\
            Example 3.2: \\
            For the same A as example 3.1, suppose Rack 2 is broken.
            \begin{enumerate}
                \item The 1-hop path [0, 7] (via rack 2) is now impossible.

                \item BFS checks other neighbors of 0. Let's say it finds node 6
                    (A[0][7] = 6, via rack 7).

                \item BFS then checks neighbors of 6. (A[6][4] = 7). This is
                    node 7, via rack 4.

                \item Rack 4 is not broken, so this path is valid. Expected: [0,
                    6, 7], Rack: [0, 7, 4] (or another valid 2-hop path)
            \end{enumerate}

        \item Weighted Indirect Path \\
            This uses Dijkstra's Algorithm to find the optimal path from node a to
            node b (not necessarily the shortest). \\
            Minimum of 1 hop and maximum of
            $\left\lceil \frac{n}{2}\right\rceil$ hops\\
            timeslots = number of hops\\
            For an arbitrary i broken racks, at least no hop up to +i hops
            change\\
            number of timeslots = maximum number of hops \\
            $\theta_{k}$ is the weight of the list where for each element, it is
            the weight of the corresponding rack (constant if equal weight) \\
            broken racks can be represented by $\infty$ weight\\
            Example 3.3: \\
            \[
                W_{8 \times 8}= \mathscr{M}\mid W_{A_{i,k},i}= \theta_{k}, \neg(W
                _{A_{i,k},i}) = 0
            \]
    \end{enumerate}

    \section{Sirius}
    $\lambda_{1}, ..., \lambda_{i},...\lambda_{j}$ = wavelength, corresponding timeslot
    \\
    $P_{1}, ..., P_{k},...P_{l}$ = ports \\
    $G_{1}, ... , G_{p} , ... G_{q}$ = gratings (AWGRs) \\
    $N_{1}, ... , N_{s} , ... N_{r}$ = nodes \\
    $r=j \times l$ (derivation is the pathways for an arbitrary node)\\
    $jq=rl$ (derivation is the total number of pathways)\\
    Then, it follows that $q = l^{2}$\\
    \
Parent Function: \\
    \[
        \mathscr{M}_{b \times n}= DW_{i} = S, i \in \{1, ... ,j\}\

    \]
    Source Matrix: $\mathscr{M}_{q \times l}\mid$
    \[
        W_{i} =
    \]
    Example 4.11 (2 wavelengths, 3 ports, 9 gratings, 6 nodes): \\
    \[
        A^{1} =
        \begin{bmatrix}
            1 & 7  & 13 \\
            4 & 10 & 16 \\
            2 & 8  & 14 \\
            5 & 11 & 17 \\
            3 & 9  & 15 \\
            6 & 12 & 18 \\
        \end{bmatrix}
        , A^{2} =
        \begin{bmatrix}
            4 & 10 & 6  \\
            1 & 7  & 13 \\
            5 & 11 & 17 \\
            2 & 8  & 14 \\
            6 & 12 & 18 \\
            3 & 9  & 15 \\
        \end{bmatrix}
    \]
    \[
        W^{1}_{8 \times 8}= \mathscr{M}\mid W_{A^1_{i,k},i}= 1, \neg(W_{A_{i,k},i}
        ) = 0, W^{2}_{8 \times 8}= \mathscr{M}\mid W_{A^1_{i,k},i}= 1, \neg(W_{A_{i,k},i}
        ) = 0
    \]
    Example 4.12 (3 wavelengths, 2 ports, 4 gratings, 6 nodes): \\
    \[
        A^{1} =
        \begin{bmatrix}
            1 & 7  \\
            3 & 9  \\
            5 & 11 \\
            2 & 8  \\
            4 & 10 \\
            6 & 12 \\
        \end{bmatrix}
        , A^{2} =
        \begin{bmatrix}
            3 & 9  \\
            5 & 11 \\
            1 & 7  \\
            4 & 10 \\
            6 & 12 \\
            2 & 8  \\
        \end{bmatrix}
        , A^{3} =
        \begin{bmatrix}
            5 & 11 \\
            1 & 7  \\
            3 & 9  \\
            6 & 12 \\
            2 & 8  \\
            4 & 10 \\
        \end{bmatrix}
    \]
\end{document}